\section{Algoritmi}
V tem razdelku bomo vzpostavili algoritmični problem izračuna negibnih točk odse\-koma linearnih funkcij in ga rešili z ustreznimi algoritmi, katerih koda bo tudi javno dostopna. Zanimali nas bosta le najmanjša in največja negibna točka, saj sta pomembni za relevantne probleme, v katerih želimo učinkovito izračunati ti vrednosti. Pojma najmanjše in največje negibne točke sta hkrati dualna: če imamo algoritem za izračun najmanjše negibne točke, lahko s pomočjo njega in ustreznih manipulacij s preslikavo izračunamo tudi največjo negibno točko. Obenem bomo zastavili vprašanja, na katera odgovor še ni znan, in poskušali poiskati ustrezne rešitve.  

\subsection{Formulacija problema}
Naj bo \(E\) polna mreža. Na množici monotonih preslikav v drugi spremenljivki \(E^{X \times E}\) definiramo operator 
\[
  \dagger_c: E^{X \times E} \to E^X, \quad f \mapsto f^{\dagger_c},
\] 
kjer je 
\[
  \all{x \in X} f^{\dagger_c}(x) = c(\Fix(y \mapsto f(x, y) \colon E \to E))
\]
in je \(c\) neka funkcija izbire elementa iz množice.

\begin{zgled}
  Naj bo \(c = \bigwedge\). Tedaj operator \(\dagger_c\) za dano monotono preslikavo v drugi spremenljivki \[f \colon X \times E \to E\] 
  vrne preslikavo 
  \[
    f^{\dagger_c} \colon X \to E,
  \] 
  ki za vsak \(x \in X\) izračuna najmanjšo negibno točko preslikave 
  \[
    f_x \colon E \to E, \quad f_x(y) = f(x, y). \qedhere
  \]
\end{zgled}

Zanimali nas bodo algoritmi, ki predstavljajo operator \(\dagger_c\) na množici odsekoma linearnih funkcij, monotonih v drugi spremenljivki (oziroma vektorjev takih funkcij), \(E^{X \times E}\), za različne izbire množic \(X\) in \(E\) ter funkcije \(c\). Oglejmo si nekaj možnih izbir.

Najenostavnejša izbira je \(X = \one\) in \(E = [0,1]\). Tako dobimo množico funkcij
\[
  [0,1]^{\one \times [0,1]}.
\]
Operator \(\dagger_c\) tedaj vrne funkcijo na množici \(\one\), pri čemer lahko vrednost \(f^{\dagger_c}\) razumemo kot negibno točko (izbrano s pomočjo funkcije \(c\)) funkcije
\[
  f \colon [0,1] \to [0,1].
\]

Prejšnji primer lahko z izbiro \(X = [0,1]^n\) posplošimo na funkcije več spremenljivk
\[
  f \colon [0,1]^n \times [0,1] \to [0,1].
\]
Operator \(\dagger_c\) tedaj vrne funkcijo na množici \([0,1]^n\):
\[
  f^{\dagger_c} \colon [0,1]^n \to [0,1].
\]
Izkaže se, da je tudi funkcija \(f^{\dagger_c}\) odsekoma linearna. Ob dodatni predpostavki, da je monotona v zadnji spremenljivki, lahko ponovno uporabimo operator \(\dagger_c\) (ki zdaj deluje na množici \([0,1]^{[0,1]^n}\)) in tako dobimo funkcijo
\[
  f^{\dagger_c \dagger_c} \colon [0,1]^{n-1} \to [0,1].
\]

Ob ustreznih predpostavkah o monotonosti lahko ta postopek nadaljujemo. Izka\-že se, da je to smiselno, vendar algoritem, predstavljen v \cite{zbMATH06805750} in posodobljen v \cite{Petkovic2017}, povzroča eksponentno časovno zahtevnost ter pri večjem številu spremenljivk hitro postane neuporaben.

Zato bomo problem poskusili rešiti z izbiro množic \(X = [0,1]^m\) in \(E = [0,1]^n\), oziroma ga bomo posplošili na vektor odsekoma linearnih funkcij
\[
  F = (f_1, \ldots, f_n) \colon [0,1]^m \times [0,1]^n \to [0,1]^n,
  \quad
  f_i \colon [0,1]^m \times [0,1]^n \to [0,1].
\]
Operator \(\dagger_c\) v tem primeru vrne preslikavo na množici \([0,1]^m\):
\[
  F^{\dagger_c} \colon [0,1]^m \to [0,1]^n.
\]

\subsection{Preizkusi časovne zahtevnosti}
Vse različice algoritma bomo testirali na skrbno izbranih primerih, zasnovanih tako, da poudarijo njihove časovne lastnosti. Pričakujemo, da bo posplošitev na vektorski primer prinesla občuten izboljšek zmogljivosti v primerjavi z iterativno uporabo skalarne različice.
